% SGA 3, Expose VI_A, Loop 1 English reconstruction.
% Scope: paragraph 3.2.3 in full; paragraph 3.2.4 is excluded.
% Authority: Polo--Gille Exp6A-13oct24.pdf, component-local pages 23--24
% (combined-reader pages 326--327; printed pages 315--316).
% OCR and the Reinhold English Markdown are comparison witnesses only.
% The commutative diagram is an exact 180-dpi crop of the controlling PDF:
% local p. 24, x=265,y=420,w=930,h=325 on a 1488x2105 render.

\medskip
\phantomsection\label{par:sga3-VIA-3.2.3}%
\noindent\textbf{3.2.3.---}
现在可以证明, 当 \(F\) 和 \(G\) 在 \(A\) 上有限型时,
\(F\backslash G\) 存在.
取在 \(A\) 上充分大的 \(A'\), 使 \(U(A')\) 与 \(G\otimes_AA'\) 一致
(参见\hyperref[par:sga3-VIA-3.2.2]{3.2.2}). 令
\[
  A''=A'\otimes_AA',
\]
且对每个 \(A\)-概形 \(X\), 用 \(X'\) 和 \(X''\)
表示纤维积 \(X\otimes_AA'\) 及 \(X\otimes_AA''\).
由\hyperref[par:sga3-VIA-3.2.1]{3.2.1} 和\hyperref[par:sga3-VIA-3.2.2]{3.2.2},
商 \(F'\backslash G'\) 和 \(F''\backslash G''\) 存在,
且有如下交换图, 其中前两行和前两列正合:
\begin{center}
  \phantomsection\label{diag:sga3-VIA-3.2.3-exact-descent}%
  \begin{tabular}{c@{\qquad}c}
    \textup{(*)}
    &
    \begin{tikzcd}[row sep=large,column sep=large]
      F''\times_{A''}G''
        \arrow[r,shift left=1.2ex,"\operatorname{pr}_2''"]
        \arrow[r,shift right=1.2ex,"\lambda''"']
        \arrow[d,shift left=1.2ex,"w_2"]
        \arrow[d,shift right=1.2ex,"w_1"']
      &
      G''
        \arrow[r,"p''"]
        \arrow[d,shift left=1.2ex,"v_2"]
        \arrow[d,shift right=1.2ex,"v_1"']
      &
      F''\backslash G''
        \arrow[d,shift left=1.2ex,"u_2"]
        \arrow[d,shift right=1.2ex,"u_1"']
      \\
      F'\times_{A'}G'
        \arrow[r,shift left=1.2ex,"\operatorname{pr}_2'"]
        \arrow[r,shift right=1.2ex,"\lambda'"']
        \arrow[d,"h"']
      &
      G'
        \arrow[r,"p'"]
        \arrow[d,"g"]
      &
      F'\backslash G'
      \\
      F\times_A G
        \arrow[r,shift left=1.2ex,"\operatorname{pr}_2"]
        \arrow[r,shift right=1.2ex,"\lambda"']
      &
      G\mathpunct{.}
      &
      {}
    \end{tikzcd}
  \end{tabular}
\end{center}

此图中, \(\operatorname{pr}'_2\) 和 \(\lambda'\)
(相应地, \(\operatorname{pr}''_2\) 和 \(\lambda''\))
由 \(\operatorname{pr}_2\) 和 \(\lambda\) 经显然的基变换得到;
态射 \(g\) 和 \(h\) 由典范单射 \(A\to A'\) 诱导.
用 \(p'\) 和 \(p''\) 表示典范态射,
而 \(v_1,v_2\) 和 \(w_1,w_2\) 由 \(A'\) 到 \(A''\) 的两个典范单射诱导.

最后, 商 \(F'\backslash G'\) 的形成与两个基变换
\[
  f_1,f_2:\operatorname{Spec}A''
    \rightrightarrows \operatorname{Spec}A'.
\]
可交换. 因此, 若
\[
  \pi':F'\backslash G'\longrightarrow\operatorname{Spec}A'
\]
表示结构态射, 则对 \(i=1,2\), 有典范同构
\[
  \tau_i:F''\backslash G''
    \xrightarrow{\ \sim\ }
    (F'\backslash G')\times_{\pi',\,f_i}\operatorname{Spec}A'',
\]
且 \(u_i\) 为 \(\tau_i\) 与投影
\[
  (F'\backslash G')\times_{\pi',\,f_i}\operatorname{Spec}A''
    \longrightarrow F'\backslash G'.
\]

的复合. 现在, 对前两行和前两列正合的 \((*)\) 型图,
容易验证 \(\operatorname{Coker}(\operatorname{pr}_2,\lambda)\) 存在
当且仅当 \(\operatorname{Coker}(u_1,u_2)\) 存在, 且此时这两个余核可等同.
因此 \(F\backslash G\) 的存在性由 \(\operatorname{Coker}(u_1,u_2)\) 的存在性得到.

\(F\backslash G\) 的形成与这里所考虑的基扩张的相容性
(参见\SGARefLink{par:sga3-VIA-3.2.1}{3.2.1} 的编者注~\ref{note:sga3-VIA-37}
及下文\SGARefLink{par:sga3-VIA-4.6}{4.6}) 表明, 复合
\begin{center}
  \phantomsection\label{diag:sga3-VIA-3.2.3-descent-datum}%
  \(
  (F'\backslash G')\times_{\pi',\,f_1}\operatorname{Spec}A''
    \xrightarrow{\ \tau_1^{-1}\ }
  F''\backslash G''
    \xrightarrow{\ \tau_2\ }
  (F'\backslash G')\times_{\pi',\,f_2}\operatorname{Spec}A''
  \)
\end{center}
为 \(F'\backslash G'\) 上关于
\[
  f:\operatorname{Spec}A'\longrightarrow\operatorname{Spec}A.
\]
的下降资料. 由\hyperref[par:sga3-VIA-3.2.1]{3.2.1}~\eqref{eq:sga3-VIA-3.2.1-dagger}
及 SGA~1, VIII.7.6, 此下降资料有效;
即 \(\operatorname{Coker}(u_1,u_2)\) 存在.
(也可直接应用\SGARefLink{thm:sga3-V-4.1}{第 V 篇, 定理~4.1}.)
